% ----------------------------------------------------------
% Fundamentação Teórica
% ----------------------------------------------------------
\chapter{REFERENCIAL TEÓRICO} \label{cha:Referencial teórico}

\section{JOGOS DIGITAIS}
O jogo digital é compreendido como um produto multidisciplinar de múltiplas camadas, que vai além da estrutura de um \textit{software} para se consolidar como uma mídia de interação e imersão \cite{silva2022elementos}.

Complementarmente, o jogo digital é caracterizado como uma atividade lúdica fundamentada em uma série de ações e decisões, circunscrita por regras e por um universo próprio que resultam em uma condição final específica \cite{alves2020jogos}. 

Além de sua dimensão lúdica e tecnológica, os jogos digitais também assumem relevância social e econômica. Sua aplicação ultrapassa o entretenimento, alcançando áreas como educação, saúde e capacitação profissional, enquanto sua produção contribui para a geração de negócios, empregos qualificados e inovações que podem alcançar outros setores da economia \cite{fortim2022pesquisa}.

\section{FEEDBACKS NO DESENVOLVIMENTO DE JOGOS DIGITAIS}

Neste trabalho, o termo \textit{feedback} é utilizado de forma ampla para abranger sugestões, ideias, avaliações e relatos relacionados aos jogos digitais.

Segundo \citeonline{sales2023feedback}, o \textit{feedback} representa uma oportunidade estratégica para expressar opiniões e sugerir caminhos para o aperfeiçoamento contínuo das atividades, permitindo o alinhamento de expectativas. 

O \textit{feedback} pode atuar como um recurso para favorecer o aprendizado contínuo e a identificação de aspectos que podem ser aprimorados em atividades e processos \cite{aquino2023feedback}. No cenário do desenvolvimento de sistemas, a estruturação desse processo é relevante para que a organização compreenda o real nível de produtividade e identifique pontos cegos que prejudicam a fluidez operacional.

Segundo \citeonline{camara2025engenharia}, a adoção de processos iterativos, testes estruturados e ferramentas de gestão da qualidade pode auxiliar na identificação antecipada de falhas e desvios, favorecendo o alinhamento entre as expectativas dos usuários, os requisitos do projeto e o produto desenvolvido.

\citeonline{pinto2024avaliacao} aplicaram um instrumento de avaliação da experiência do jogador em três jogos desenvolvidos por estudantes. Os resultados permitiram reunir percepções sobre aspectos técnicos, controles, equilíbrio, elementos visuais e sonoros, demonstrando como avaliações estruturadas podem contribuir para identificar pontos fortes e aspectos que podem ser aprimorados.

\section{CENTRALIZAÇÃO E ORGANIZAÇÃO DE FEEDBACKS}

As avaliações publicadas por jogadores constituem uma fonte relevante para compreender suas preferências e níveis de satisfação. Ao analisarem 854.505 avaliações da plataforma Steam, Amaral e Cavalcanti identificaram que o tratamento estruturado desses textos permite reconhecer padrões de sentimento, engajamento e afinidade entre comunidades, demonstrando o valor da organização dos \textit{feedbacks} em larga escala \cite{amaral2024desvendando}.

A centralização de \textit{feedbacks} em plataformas web permite estruturar avaliações, comentários e sugestões em entidades vinculadas a jogos específicos. No trabalho de Souza, essa organização é realizada por meio de avaliações, categorias, listas e tópicos temáticos, possibilitando que informações produzidas pelos usuários sejam armazenadas, consultadas e relacionadas em um ambiente integrado \cite{souza2024ludus}.

\citeonline{esteves2025uso} observa que a fragmentação de informações em diferentes plataformas pode dificultar o acesso, a comparação e a análise de dados relevantes para a tomada de decisão. No estudo realizado, a centralização dos processos exigiu a unificação das informações e a padronização de indicadores, enquanto relatórios e painéis de controle passaram a oferecer uma visão mais ampla do desempenho organizacional.

\citeonline{de2026fragmentaccao} identificam que a dispersão de informações entre sistemas, planilhas e documentos distintos dificulta o acesso, a consolidação e a verificação dos dados, além de favorecer redundâncias, retrabalho e limitações de rastreabilidade. Nesse sentido, a centralização e a padronização das informações apresentam-se como medidas relevantes para reduzir controles paralelos e ampliar a confiabilidade, a consulta e o aproveitamento dos registros produzidos.

\section{GESTÃO DE DADOS EM APLICAÇÕES WEB}

\citeonline{santos2023representaccao} identifica que a organização das informações sobre jogos na plataforma Steam é prejudicada pela ausência de relações hierárquicas bem definidas entre os marcadores utilizados na classificação dos conteúdos. Como alternativa, o estudo propõe a integração entre um vocabulário controlado e a participação colaborativa dos usuários, buscando tornar a representação, a categorização e a recuperação das informações mais consistentes dentro da plataforma.

Além da organização e da recuperação dos dados, plataformas que recebem conteúdos produzidos por usuários também precisam considerar aspectos relacionados à moderação e à governança dessas contribuições.

\citeonline{ribeiro2025engajamento} apontam que o conteúdo gerado pelos usuários amplia a interação e valoriza sua participação em plataformas digitais, ao mesmo tempo em que exige mecanismos de controle para reduzir riscos relacionados à desinformação, ao uso indevido e à confiabilidade das informações compartilhadas. Nesse sentido, ambientes que recebem contribuições dos usuários precisam combinar incentivo à participação com formas de organização, moderação e governança dos conteúdos produzidos.

\citeonline{melo2020descoberta} demonstram que conversas geradas em chats de jogos online podem ser tratadas como uma fonte de dados capaz de revelar padrões e associações relevantes sobre o comportamento e as estratégias dos jogadores. No estudo realizado com o jogo Mobile Strike, a aplicação conjunta das ferramentas \textit{Text Mining Suite} e Sobek possibilitou extrair novos conceitos e gerar regras associativas com maior frequência e qualidade, evidenciando o potencial da mineração de textos para organizar e analisar informações não estruturadas produzidas durante as interações dos usuários.
